\documentclass[11pt,a4paper]{article}
\usepackage[utf8]{inputenc}
% \usepackage[T1]{fontenc}  % disabled: no scalable T1 fonts on the build machine (would embed bitmap Type 3 fonts)
\usepackage{amsmath,amssymb,amsthm}
\DeclareTextSymbolDefault{\textbullet}{OMS}            % take the two symbols that would otherwise come from the
\DeclareTextSymbolDefault{\textasteriskcentered}{OMS}  % bitmap-only TS1 fonts from the Type 1 cmsy font instead
\usepackage{booktabs}
\usepackage{array}
\usepackage[margin=2.6cm]{geometry}
\usepackage[expansion=false]{microtype}
\usepackage[hidelinks]{hyperref}
\usepackage{url}

\newtheorem{lemma}{Lemma}
\newtheorem{observation}{Observation}
\theoremstyle{remark}\newtheorem{remark}{Remark}

\newcommand{\F}{\mathbb{F}}
\newcommand{\B}{\mathcal{B}}
\newcommand{\Dreg}{D_{\mathrm{reg}}}
\newcommand{\Dsolve}{D_{\mathrm{solve}}}
\newcommand{\lm}{\mathrm{LM}}
\newcommand{\std}{\mathrm{std}}
\newcommand{\Tr}{\mathrm{Tr}}
\newcommand{\SR}{\mathrm{SR}}

\title{The first fall degree assumption and the degree of regularity of Weil-descent systems for summation polynomials over $\F_{2^n}$:\\
exact closure dimensions, thresholds and a growth law}
\author{Alexander Solovyev\thanks{Independent researcher, Omsk, Russia. E-mail: \texttt{alekssolo118@gmail.com}. ORCID: 0009-0003-3657-863X.}}
\date{September 2026}

\begin{document}
\maketitle

\begin{abstract}
The only known route to a subexponential-time solution of the discrete logarithm problem on elliptic curves over binary fields $E(\F_{2^n})$ is the index calculus of Gaudry and Diem, with relations found by Weil descent of Semaev's summation polynomials. Its heuristic running time rests on the \emph{first fall degree assumption} of Petit and Quisquater: that the degree of regularity of the descended systems stays close to their first fall degree as $n$ grows. Huang, Kosters and Yeo (CRYPTO 2015) tested this for the third summation polynomial $S_3$ up to $n=40$ in Magma, saw the degree rise from 3 to at least 5, and left its growth open. We settle the question for $S_3$ as far as computation allows and take the first measurements for $S_4$. Three independent tools are used: an exact computation of the degree-capped closures $V_D$ of their Definition~1 (32 variables at degree 5; 24 variables at degree 7), a corrected build of the M4GB Gröbner-basis solver, and a Buchberger reference implementation. For $S_3$ (two-point decompositions): (i)~the degree of regularity is a deterministic function of the pair $(N,n)$ of variables and equations; (ii)~the generic closures have exact polynomial dimensions, $\dim V_3=(3N-2)n-\binom N2$, $\dim V_4=\tfrac12[3N(N+1)n-171n-9n^2]$ and a degree-5 part $Q_5=3n\binom N3-\binom{3n+1}2N-n(79N-660)$, verified to the unit on 15 further systems; (iii)~degree $D$ suffices exactly when the generic degree-$D$ part exceeds $\binom ND$, a criterion validated on 59 boundary cases and 110 grid outcomes without exception; (iv)~the ideal contains exactly $3n$ quadrics and the closures behave like a semi-regular system of $3n$ Boolean quadrics up to a small lag. Consequently the degree of regularity of $S_3$ is 3 for $n\le16$, 4 for $17\le n\le36$, 5 for $37\le n\le52$, and grows linearly (slope $\approx1/20$ here, $0.036$ asymptotically): the first fall degree assumption fails for two-point decompositions, quantitatively. For $S_4$ (three-point decompositions, degree-6 equations) the computed range shows the opposite: in all twenty systems computed, from square systems with 12 to 24 variables down to systems with three times as many variables as equations, the degree-7 closure contains the whole truncated ideal, so the solving degree equals the first fall degree 7 and no generic regime appears at all. We also report and fix a deviation from the Gebauer--Möller criteria in M4GB (merged upstream), and find that M4GB is unsuited to the dense degree-6 systems of $S_4$, for which the closure computation is the practical tool.
\end{abstract}

\paragraph{Notation.} $N$ is the number of Boolean variables and $n$ the number of quadratic equations (the extension degree); $\B=\F_2[x_1,\dots,x_N]/(x_i^2+x_i)$ is the Boolean polynomial ring, $\B_{\le k}$ its subspace of polynomials of degree $\le k$, $|\B_{\le k}|=\sum_{j\le k}\binom Nj$; $\binom{N}{\le k}$ abbreviates the same sum.

\section{Introduction}

For an elliptic curve over a prime field no algorithm for the discrete logarithm problem is known that beats the generic $\sqrt{\#E}$ bound. Over extension fields the situation is different: the index calculus of Gaudry \cite{Gaudry09} and Diem \cite{Diem11,Diem13} finds relations by \emph{point decomposition} --- given $Q$, find $P_1,\dots,P_m$ in a factor base $\mathcal F_V=\{P: x(P)\in V\}$, $V$ an $\F_q$-subspace of $\F_{q^n}$, with $P_1+\dots+P_m=Q$ --- and turns each decomposition into a polynomial system by Weil descent of Semaev's summation polynomial $S_{m+1}$ \cite{Semaev04}. For binary fields Petit and Quisquater \cite{PQ12} observed that the descended systems have a very low \emph{first fall degree} and argued that, if the degree at which a Gröbner-basis computation terminates stays close to it as $n$ grows --- the \emph{first fall degree assumption} --- the whole attack runs in time $2^{O(n^{2/3}\log n)}$, subexponential in the field size. Whether the assumption holds is therefore the question on which the asymptotic security of elliptic curves over $\F_{2^n}$ against this attack turns. Huang, Kosters and Yeo \cite{HKY15} gave a rigorous treatment of the related \emph{last fall degree} and tested the assumption experimentally in the two-point case ($S_3$, with $x(P_0),x(P_1)$ in a random subspace of dimension $\lceil n/2\rceil$): in Magma, for $n\le40$, the degree of regularity rose from 3 to at least 5 without a discernible law, and they posed its growth as an open problem \cite[Sec.~5.2]{HKY15}; Kosters and Yeo \cite{KY15} added $n=17$ and $30$ and reported that $n=40$ exhausted 38\,GB of memory before finishing.

This note takes that question as far as computation allows, for $S_3$ and, for the first time, for $S_4$, and then one step further by a structural model. Our contributions are the following.

\begin{enumerate}
\item \textbf{Deterministic dependence on $(N,n)$.} For random subspaces the degree of regularity depends only on the number of variables $N$ and equations $n$: it is independent of the curve, the subspace, the point $Q$, the number of solutions (including inconsistent systems), and of the choice of instance. Under the convention $\dim V=\lceil n/2\rceil$ of HKY, odd $n$ has one variable more than equations ($N=n+1$) and reaches the next degree earlier than even $n$ --- the source of the apparently irregular entries $15\to4$, $35\to5$ in our tables.

\item \textbf{Exact closure dimensions.} Let $V_D$ be the degree-capped closure of the system (HKY, Definition~1). In the \emph{generic} regime --- before the closure collapses onto the whole truncated ideal --- its dimensions are exact polynomials in $(N,n)$:
\[
\dim V_3=(3N-2)n-\binom N2,\qquad \dim V_4=\tfrac12\bigl[3N(N+1)n-171n-9n^2\bigr],
\]
and the degree-5 part of $V_5$ is $Q_5=3n\binom N3-\binom{3n+1}2N-n(79N-660)$. These were fitted on a handful of points and then verified to the unit on 9 further generic points at $N\le32$ ($\dim V_5$ up to $221\,331$) and on 6 collapsed points, with margins down to $0.15\%$.

\item \textbf{A threshold criterion.} Degree $D$ suffices --- the reduced Gröbner basis lies in $V_D$, so no F4-type algorithm needs a higher step degree --- exactly when the generic degree-$D$ part exceeds the number of degree-$D$ monomials: $Q_D(N,n)>\binom ND$. The criterion was tested on 59 boundary cases and 110 further grid outcomes (square, odd and unbalanced systems; closure computations and two Gröbner-basis solvers) without exception. It gives $\Dreg=3$ for $n\le16$, 4 for $17\le n\le36$ and 5 for $37\le n\le52$, the last value being a prediction; direct computations reach $n=40$ (M4GB, 6\,h) and confirm the value 5 that Magma could not reach \cite{KY15}.

\item \textbf{Structure and a growth law.} The ideal contains exactly $3n$ quadratic polynomials: the $n$ equations, $N-1$ multiples of the trace-linear form of HKY's Proposition~5, and $2n-N+1$ further quadrics arising as degree falls. The closures behave, up to an explicit lag, like the Macaulay matrices of a \emph{semi-regular} system of $3n$ Boolean quadrics in $N$ variables \cite{BFS03,BFSY05}: the semi-regular thresholds are $20,37,56,76,97,119,\dots$ for $D=3,4,5,\dots$, ours are $16,36,52$. Consequently $\Dreg$ grows linearly in $n$, with about one degree per 20 variables in the computed range and slope $0.036$ asymptotically. \emph{The first fall degree assumption fails for two-point decompositions}, and it fails in a quantified way: the systems are exactly as easy as generic systems with three times as many equations, and no easier --- a constant factor in the exponent relative to generic systems of $n$ quadrics (slope $0.09$), not a change of complexity class.

\item \textbf{Three-point decompositions.} The Weil descent of $S_4$ yields $n$ equations of degree 6 in $N=3\lceil n/3\rceil$ variables. With a memory-light closure implementation reaching 24 variables at degree 7 we find the opposite behaviour in the computed range: in all twenty systems computed --- square systems with $N=12,\dots,24$ and systems with as few as $N/3$ equations --- the degree-7 closure contains the whole truncated ideal, so the solving degree equals the first fall degree 7 and no generic regime appears, whereas for $S_3$ the generic regime sets in at $N=17$ and its boundary then moves by about $0.025N$ per variable. The assumption of Petit and Quisquater is thus consistent with all data for $m=3$ while refuted for $m=2$.

\item \textbf{Solvers.} The reported step degree of the Gröbner-basis solver M4GB \cite{MS17} was inflated by one on part of the $S_3$ family because its pair update discarded both members of every pair of critical pairs with equal lcm, whereas the Gebauer--Möller criteria keep one. We located, reproduced and fixed the defect; the fix was merged upstream (pull request \#14, September 2026), and the corrected solver agrees with the closure computations everywhere. On the dense degree-6 systems of $S_4$, M4GB exhausts memory already at 21 variables (matrices of $2\cdot10^6$ rows), so that there the closure computation, not the Gröbner-basis solver, is the practical instrument.
\end{enumerate}

We stress what is \emph{not} claimed. The formulas are empirical: exact on every point we could compute, with the structural explanation of Section~\ref{sec:structure} reaching the first of the two stages of the closure. Beyond $D=5$ the $S_3$ thresholds are predictions of the semi-regular model with a lag that we can bound but not derive. For $S_4$ the computed range is $N\le24$, and twenty collapses are evidence, not proof, that the assumption holds; the question is settled only for $m=2$. No practical consequence follows for curves in use: two-point decompositions are exponential regardless, and for $m\ge3$ the cost of a single decomposition at cryptographic sizes remains far beyond anything reached here.

\section{Preliminaries}

\subsection{Curves, summation polynomial, Weil descent}

Let $k=\F_{2^n}$ and let $E: y^2+xy=x^3+a_2x^2+a_6$ be an ordinary elliptic curve over $k$ (every ordinary curve over $k$ is isomorphic to one of this form; the isomorphism acts on $x$ by an affine map and does not affect the degrees considered below). The third summation polynomial of $E$ in characteristic 2 is, from the general formula of \cite[Sec.~5.1]{HKY15} with $b_2=1$, $b_4=b_6=0$, $b_8=a_6$,
\[
S_3(X_0,X_1,X_2)=X_0^2X_1^2+X_0^2X_2^2+X_1^2X_2^2+X_0X_1X_2+a_6,
\]
with the property that $S_3(x(P_0),x(P_1),x(P_2))=0$ if and only if $P_0\pm P_1\pm P_2=O$ for some choice of signs; we verified this on random triples for every field used.

Fix a subspace $V\subset k$ of dimension $d=\lceil n/2\rceil$ with basis $v_1,\dots,v_d$ and a point $Q\in E(k)$. Substituting $X_0=\sum u_iv_i$, $X_1=\sum w_jv_j$ with $u_i,w_j\in\F_2$ and expanding $S_3(X_0,X_1,x(Q))$ in a basis of $k$ over $\F_2$ gives $n$ polynomials in $N=2d$ variables. Because squaring is $\F_2$-linear, $X_0^2X_1^2$ and $X_0^2x(Q)^2$ contribute only bilinear and linear monomials: the system is
\[
f_r(u,w)=\sum_{i,j}c_{r,ij}\,u_iw_j+\sum_i d_{r,i}\,u_i+\sum_j e_{r,j}\,w_j+a_{6,r},\qquad r=1,\dots,n,
\]
with $c_{ij}=v_i^2v_j^2+v_iv_jx(Q)$, $d_i=v_i^2x(Q)^2$, $e_j=v_j^2x(Q)^2$ (coordinates over $\F_2$). Together with the field equations $u_i^2=u_i$, $w_j^2=w_j$ this is the system studied in \cite[Sec.~5.2]{HKY15}; we work in the Boolean ring $\B$, where the field equations are built in. The solution set is the set of decompositions of $Q$; by the symmetry $X_0\leftrightarrow X_1$ it has even size unless $X_0=X_1$ is possible. For even $n$ the system is square ($N=n$); for odd $n$ it has one variable more than equations ($N=n+1$).

HKY take $Q$ to be a random point, so that about 60\% of their systems have no solution; we generate both random-$Q$ systems and \emph{consistent} systems $Q=P_0+P_1$ with $P_0,P_1$ drawn uniformly from $\mathcal F_V$, and find no difference in any degree measured. We also consider unbalanced systems in which $V$ has dimension $d>\lceil n/2\rceil$ or the two points use subspaces of different dimensions $d_0+d_1=N$; the family is parametrised by $(N,n)$ with $n\ge\max(d_0,d_1)$, hence $N\le 2n$.

\subsection{Degree-capped closures and degrees}

For a set $F$ of polynomials in $\B$ and $D\ge2$, the closure $V_D=V_{F,D}$ is the smallest $\F_2$-subspace containing every $f\in F$ of degree $\le D$ and closed under $g\mapsto x_a g$ whenever $\deg g\le D-1$ (\cite[Definition~1]{HKY15}, transported to the Boolean ring; the two definitions agree modulo the field equations). We call the \emph{first fall degree} the least $c$ with $V_c\cap\B_{\le c-1}\ne V_{c-1}$, the \emph{last fall degree} the greatest such $c$, and the \emph{solving degree} $\Dsolve$ the least $D$ such that $V_D$ contains the reduced Gröbner basis of the ideal (for a degree-compatible order; cf.\ \cite{CG23,CCG23}). For the systems at hand the first fall degree is always 2 (Proposition~5 of \cite{HKY15}: a fixed linear combination of the $f_r$ is the linear polynomial $\Tr((x(Q)+a_2)/a_1^2)+\Tr(X_0+X_1)=0$), and the last fall degree coincided with $\Dsolve$ in every instance we computed.

HKY define the degree of regularity operationally as the largest step degree of Magma's F4 at which new polynomials are produced. For any F4-type algorithm with the normal selection strategy, every polynomial produced at step degree $\le D$ lies in $V_D$; hence

\begin{lemma}\label{lem:dsolve}
$\Dsolve\le\Dreg$ for every such algorithm.
\end{lemma}

The inequality is an equality for Magma's F4 on all of HKY's data points and for the corrected M4GB on all of ours; the uncorrected M4GB violates it by one on part of the family (Section~\ref{sec:m4gb}). We therefore use $\Dsolve$ --- a property of the system, not of a program --- as the primary quantity and report the solver's degree as a check.

\subsection{Testing containment of the Gröbner basis}

Let $E_D$ be an echelon basis of $V_D$ with respect to a monomial order refining degree (we use grevlex) and let $\lm(V_D)$ be its set of leading monomials. Write $\std(S)$ for the number of squarefree monomials divisible by no element of a set $S$.

\begin{lemma}\label{lem:std}
$V_D$ contains the reduced Gröbner basis of the ideal $J$ generated by $F$ if and only if $\std(\lm(V_D))$ equals the number of solutions of $F$ in $\F_2^N$.
\end{lemma}

\begin{proof}
$J$ is radical and $\B/J$ has dimension $\#V(F)$, so $\std(\lm(J))=\#V(F)$; $\lm(V_D)\subseteq\lm(J)$ gives $\std(\lm(V_D))\ge\#V(F)$, with equality iff the two monomial ideals coincide. If they coincide, every minimal generator $m$ of $\lm(J)$ is a leading monomial of $V_D$; every non-standard monomial of degree $\le D$ is then also in $\lm(V_D)$ (multiply the element of $V_D$ whose leading monomial divides it by the complementary squarefree variables --- the degree bound is respected and, in the Boolean ring, the product's leading monomial is the intended one), so the fully reduced row of $E_D$ with pivot $m$ has only standard monomials in its tail and equals the reduced Gröbner-basis element with that leading monomial.
\end{proof}

The number of solutions is obtained exactly by exploiting the bilinear structure: for each of the $2^{d_0}$ values of $u$ the system is affine in $w$ and is solved by linear algebra.

\section{Tools and methodology}

\subsection{Instance generation}

Field arithmetic over $\F_{2^n}$ (irreducible polynomials found by the Ben-Or test), curve arithmetic and the derivation of $S_3$ were implemented from scratch and checked against the group law on random points; the Weil descent was checked by evaluating every generated system on its known solution. Subspaces are random (a uniformly chosen independent set); coordinates of points in $V$ are computed by Gaussian elimination. All instances are reproducible from a seed; the seeds of every instance in this note are listed in the supplementary archive.

\subsection{Closure computations}\label{sec:tools}

Closures were computed by two independent programs. The first, in Python, represents Boolean polynomials as bit vectors over the monomials of degree $\le D$ in grevlex order and computes $V_D$ by iterated multiplication with echelon reduction; it handles $N\le26$ at $D=4$ and was used for all $D=3,4$ data and for the verification of Lemma~\ref{lem:std} by comparison with the Buchberger--Möller basis of the vanishing ideal of the solution set. The second, in C (about 400 lines), stores rows as packed 64-bit words up to their pivot, keeps the echelon block-triangular (every pivot row has zeros at the pivot columns of earlier blocks; the rows of one block are in reduced echelon form among themselves), and reduces batches of new rows by Gray-code tables of $k=10$ pivot rows (the Method of Four Russians), so that no reduced-row-echelon maintenance is needed when pivots are appended. Rows are expanded once, in insertion order, and only if their degree is $\le D-1$; with OpenMP the reduction of a batch is parallel over rows. It reproduces every Python result to the unit and reaches $N=32$ at $D=5$ (243\,000 columns, 3.7\,GB, 5--13 hours on eight threads). A second mode computes the rank of the plain product set $\{t\,q: q\in W_2,\ \deg t\le D-2\}$ of Section~\ref{sec:structure}, without the iterated expansion of degree falls. A third implementation, written for the $S_4$ systems of Section~\ref{sec:s4}, keeps the echelon in \emph{fully} reduced row echelon form and stores each row only over the columns that are not yet pivots, compacting lazily; since the rows shrink as pivots accumulate, both memory and work decrease as the closure fills the truncated ideal, and 24 variables at degree 7 ($536\,000$ columns) fit in 10\,GB where the block-triangular echelon would need 18\,GB. All three implementations agree to the unit on every system computed with more than one of them.

\subsection{Gröbner-basis solvers}\label{sec:m4gb}

We used M4GB \cite{MS17} (GF(2), field equations included as input polynomials; Cygwin, GCC~13, \texttt{-O3}) as the Gröbner-basis solver, reading $\Dreg$ from its per-step output as the largest lcm degree of a step reporting new polynomials --- the definition of HKY. Two remarks are necessary for reproducibility. First, the libtool build places a 27\,KB wrapper in \texttt{bin/} and the actual executable in \texttt{bin/.libs/}; renaming the wrapper does not preserve a build, and the executable in \texttt{.libs/} must be copied after each build with a different \texttt{MAXVARS}. Second, in the pair update of \texttt{m4gb.hpp} the Gebauer--Möller M-criterion is implemented with non-strict divisibility of lcms,
\begin{verbatim}
if (it != it2 && it2->first.intlcm | it->first.intlcm)
    { it->second = false; break; }
\end{verbatim}
so that two admissible new pairs with equal lcm discard each other, whereas the F-criterion keeps exactly one (none if a coprime pair shares the lcm) \cite{GM88}. The final basis was correct in every run we made --- the lost S-polynomials reappear through pairs of higher lcm degree --- but the maximal step degree is inflated by one where those S-polynomials carry the degree falls. On the $n=16$ systems of HKY the solver reports 4 (13 new polynomials at degree 4: exactly the 13 linear polynomials that complete the basis) against HKY's 3; with the criterion corrected (a 20-line change, submitted as pull request \#14 and merged into the main M4GB repository on 25 September 2026, \url{https://github.com/cr-marcstevens/m4gb/pull/14}; also provided as a standalone patch in the supplementary archive) it reports 3 on all 41 instances we tested, with the same basis, and agrees with $\Dsolve$ on every instance of this note. We reproduced the mechanism independently with a Buchberger implementation whose pair criteria can be switched between the two versions.

\subsection{Verification chain}

Each number in Sections~\ref{sec:results}--\ref{sec:criterion} rests on at least two of: the closure computation ($\Dsolve$ by Lemma~\ref{lem:std}), the corrected M4GB ($\Dreg$ by HKY's definition), the Buchberger reference implementation, and, for HKY's own sizes, their published Magma values. Where a polynomial formula is stated, the points used to fit it are separated from the points used to verify it, and the latter --- 21 in total --- were computed after the prediction was recorded. The threshold predictions of Section~\ref{sec:criterion} were likewise written down, with their margins, before the corresponding systems were run.
\section{Results for the systems of HKY}\label{sec:results}

Table~\ref{tab:square} collects the degree of regularity for the systems of \cite[Sec.~5.2]{HKY15} --- random curve, random subspace of dimension $\lceil n/2\rceil$ --- from $n=12$ to $n=40$. ``Closure'' is $\Dsolve$ computed by Lemma~\ref{lem:std}; ``M4GB'' is the corrected solver's largest productive step degree; the last two columns are the published values of HKY (Magma F4) and KY. For $n\le24$ the closure column summarises 378 instances (random and consistent $Q$, fixed and varying curves and subspaces, several instances per $n$) with no variation whatsoever; for $n\ge26$ two instances were run at $26$--$32$, $35$ and $36$ and one otherwise. For $n\ge26$ the closure column was computed with the memory-light implementation of Section~\ref{sec:tools}: the degree-4 closures collapse for $n=26$--$34$ and $36$ and do not collapse for $n=35$ and $37$--$40$ (so that $\Dsolve\ge5$ there), and the degree-5 closures collapse for $n=35$ and $37$--$40$: $\dim V_5=|\B_{\le5}|-\#\text{sol}$ with standard monomials of degree $\le1$ only. The $n=40$ system, which Magma could not finish in 38\,GB \cite{KY15}, took 32 hours and a peak of 20\,GB (partly in swap) on a desktop.

\begin{table}[ht]\centering\small
\caption{Degree of regularity of the Weil-descent system of $S_3$, $\dim V=\lceil n/2\rceil$. Bold: odd $n$ at which the degree increases one step earlier than for the neighbouring even $n$.}\label{tab:square}
\begin{tabular}{rrcccccl}
\toprule
$n$ & $N$ & closure & M4GB (fixed) & HKY & KY & \multicolumn{2}{l}{time / memory (M4GB)}\\
\midrule
12 & 12 & 3 & 3 & 3 & 3 & $<1$\,s\\
13 & 14 & 3 & --- & & & \\
14 & 14 & 3 & 3 & & & $<1$\,s\\
15 & 16 & \textbf{4} & --- & & & \\
16 & 16 & 3 & 3 (4 before the fix) & 3 & 3 & $<1$\,s\\
17 & 18 & 4 & --- & & 4 & \\
18 & 18 & 4 & 4 & 4 & & 1\,s\\
19 & 20 & 4 & --- & & & \\
20 & 20 & 4 & 4 & 4 & 4 & 3\,s\\
22 & 22 & 4 & --- & & & \\
24 & 24 & 4 & --- & 4 & & \\
26 & 26 & 4 & 4 & & & 26\,s / 7\,MiB\\
28 & 28 & 4 & 4 & & & 71\,s / 7\,MiB\\
30 & 30 & 4 & 4 & & 4 & 214\,s / 7\,MiB\\
32 & 32 & 4 & 4 & & & 8\,min / 9\,MiB\\
33 & 34 & 4 & 4 & & & 18\,min / 0.6\,GiB\\
34 & 34 & 4 & 4 & & & 17\,min / 0.6\,GiB\\
35 & 36 & \textbf{5} & \textbf{5} & & & 36--44\,min / 0.6--1.0\,GiB\\
36 & 36 & 4 & 4 & & & 38--39\,min / 0.5--0.8\,GiB\\
37 & 38 & 5 & 5 & & & 82\,min / 1.0\,GiB\\
38 & 38 & 5 & 5 & & & 78\,min / 0.8\,GiB\\
39 & 40 & 5 & 5 & & & 5.1\,h / 2.2\,GiB\\
40 & 40 & 5 & 5 & & $\ge5$ & 6.0\,h / 1.5\,GiB\\
\bottomrule
\end{tabular}
\end{table}

Three observations organise the rest of the paper. (i)~The degree is a function of $(N,n)$ only: 378 instances at $n\le24$ --- solvable and unsolvable, with 0 to 8 solutions, with the same or different curves and subspaces --- give identical closure dimensions in every degree, not merely the same $\Dsolve$ (Section~\ref{sec:generic}). (ii)~With $N=n+1$ for odd $n$, the thresholds are crossed one step earlier: $15\to4$ while $16\to3$, and $35\to5$ while $36\to4$; the natural coordinates are $(N,n)$, and in them there is no irregularity. (iii)~The values 3, 4, 5 are reached at $n=12$--$16$, $17$--$36$, $37$--$40$ respectively, consistent with HKY and KY at every common size and extending KY's ``$\ge5$'' at $n=40$ to an exact 5.

The unsolvable systems of HKY's construction (random $Q$) deserve a remark: the closure then contains 1, and the criterion of Section~\ref{sec:criterion} applies verbatim --- the ``degree of regularity'' of an inconsistent system is the degree at which 1 is reached, and it coincides with the degree of the consistent systems of the same $(N,n)$ in all 60 cases we compared.

\section{Generic closure dimensions}\label{sec:generic}

\subsection{Profiles}

For every instance we computed the full profile $\dim(V_D\cap\B_{\le k})$, $k\le D$. In the \emph{generic} regime --- as long as $V_D$ is a proper subspace of the truncated ideal $J_{\le D}$ --- the profiles are the same for all instances with the same $(N,n)$, and they are polynomials in $(N,n)$. In the \emph{collapsed} regime $V_D=J_{\le D}$, of dimension $|\B_{\le D}|$ minus the number of solutions (all standard monomials of these ideals have degree $\le2$), and $\Dsolve\le D$.

The lowest levels are uniform across the family:
\begin{itemize}
\item $\dim V_2=n+N-1$: the $n$ equations plus the $N-1$ independent multiples $x_a\ell$ of the trace-linear polynomial $\ell\in\operatorname{span}(f_r)$.
\item $\dim(V_3\cap\B_{\le2})=3n$: the space $W_2$ of \emph{all} quadratic polynomials of the ideal (it equals $J_{\le2}$, as confirmed by Buchberger--Möller on the solution sets) has dimension exactly $3n$. It consists of the $n$ equations, the $N-1$ multiples of $\ell$, and $M_3=2n-N+1$ further quadrics obtained as degree falls from cubic products. No further quadric appears at any higher level.
\item $\dim(V_4\cap\B_{\le3})=3n(N+1)$, and $\dim(V_5\cap\B_{\le4})=\dim V_4+8n$.
\end{itemize}

\subsection{The polynomials}

\begin{table}[ht]\centering\footnotesize
\caption{Generic closure dimensions. $P_D=\dim V_D$; $Q_D=\dim V_D-\dim(V_D\cap\B_{\le D-1})$ is the degree-$D$ part.}\label{tab:formulas}
\begin{tabular}{llp{3.6cm}p{4.6cm}}
\toprule
level & formula & fitted on & verified on\\
\midrule
$P_3$ & $(3N-2)n-\binom N2$ & $N=16,18$ (slope and intercept in $n$) & 40 further $(N,n)$, $N=14$--$26$, exact\\
$P_4$ & $\frac12[3N(N+1)n-171n-9n^2]$ & $N=18,22,24$ & $N=20,26,28,30,32$ (13 points) and the square and odd systems at $N=36,38,40$ (5 points), exact\\
$Q_5$ & $3n\binom N3-\binom{3n+1}2N-n(79N-660)$ & $(24,12)$, $(26,13)$ & 7 points, $N=28$--$32$, $n-N/2\le3$, exact\\
\bottomrule
\end{tabular}
\end{table}

The data behind Table~\ref{tab:formulas}: for $D\le4$ a grid of 55 systems with $N=16,18,\dots,24$ and $n$ running from $N/2$ to $N$ ($\dim V=N/2$ for both points, $n$ equations from the field $\F_{2^n}$); a second family with $\dim V_0=(n+1)/2$, $\dim V_1=(n-1)/2$ ($N=n$ for odd $n$); and the square and odd systems of Table~\ref{tab:square} up to $N=26$; for $P_4$ also the five generic systems $(36,35)$, $(38,37)$, $(38,38)$, $(40,39)$, $(40,40)$, whose degree-4 closures ($61\,425$, $72\,927$, $74\,727$, $85\,761$, $87\,780$) agree with the formula to the unit. For $D=5$, the 15 systems of Table~\ref{tab:level5}.

\begin{table}[ht]\centering\small
\caption{Level 5 (C program; dimensions to the unit). ``Full'' means $\dim V_5=|\B_{\le5}|-\#\text{solutions}$. The first two rows fixed the two free parameters of $Q_5$ (the coefficients 79 and 660); every other entry was predicted before it was computed.}\label{tab:level5}
\begin{tabular}{lrrrrrr}
\toprule
$(N,n)$ & $\dim(V_5\cap\B_{\le4})$ & $P_4+8n$ & $\dim V_5$ & $Q_5$ measured & $Q_5$ formula & $\binom N5$\\
\midrule
(24,12) & 9\,222 & 9\,222 & 51\,270 & 42\,048 & 42\,048 & 42\,504\\
(24,13) & 10\,909 & --- & 53\,413 full & --- & 44\,148 & 42\,504\\
(26,13) & 11\,921 & 11\,921 & 74\,919 & 62\,998 & 62\,998 & 65\,780\\
(26,14) & --- & --- & 79\,556 full & --- & 66\,206 & 65\,780\\
(26,15) & --- & --- & 81\,565 full & --- & 69\,180 & 65\,780\\
(28,14) & 15\,085 & 15\,085 & 105\,665 & 90\,580 & 90\,580 & 98\,280\\
(28,15) & 16\,095 & 16\,095 & 111\,255 & 95\,160 & 95\,160 & 98\,280\\
(28,16) & --- & --- & 118\,274 full & --- & 99\,488 & 98\,280\\
(30,15) & 18\,750 & 18\,750 & 144\,750 & 126\,000 & 126\,000 & 142\,506\\
(30,16) & 19\,928 & 19\,928 & 152\,168 & 132\,240 & 132\,240 & 142\,506\\
(30,17) & 21\,097 & 21\,097 & 159\,307 & 138\,210 & 138\,210 & 142\,506\\
(30,18) & --- & --- & 170\,182 full & --- & 143\,910 & 142\,506\\
(32,16) & 22\,952 & 22\,952 & 193\,512 & 170\,560 & 170\,560 & 201\,376\\
(32,19) & 26\,999 & 26\,999 & 221\,331 & 194\,332 & 194\,332 & 201\,376\\
(32,20) & 37\,423 & --- & 238\,799 full & --- & 201\,680 & 201\,376\\
\bottomrule
\end{tabular}
\end{table}

A second instance of $(28,15)$ and of $(28,16)$ gave the same numbers. The $8n$ new quartic polynomials at level 5 are visible in the second and third columns of Table~\ref{tab:level5}.

\subsection{What the closure is not}

The closure is not the Macaulay matrix. The rank $R_D$ of the plain product set $\{t\,q: q\in W_2,\ \deg t\le D-2\}$ is strictly smaller than $P_D$: at $(20,10)$, $R_4=2994$ against $P_4=4995$. The difference is contributed by products of variables with \emph{degree falls} --- polynomials of degree $<D$ obtained as combinations of products whose top-degree parts cancel --- which the iterated closure multiplies again. This two-stage structure is analysed in Section~\ref{sec:structure}; it is the reason why $P_D$ exceeds any single Macaulay-type count and why ``semi-regularity'' enters only after the falls are accounted for.

\section{The threshold criterion}\label{sec:criterion}

\subsection{Statement}

\begin{observation}[criterion]\label{obs:criterion}
For $D=3,4,5$ put
\begin{align*}
Q_3(N,n)&=3nN-5n-\tbinom N2,\\
Q_4(N,n)&=3n\tbinom N2-\tbinom{3n+1}2-87n,\\
Q_5(N,n)&=3n\tbinom N3-\tbinom{3n+1}2N-n(79N-660).
\end{align*}
Then $\Dsolve\le D$ if and only if $Q_D(N,n)>\binom ND$.
\end{observation}

The ``only if'' direction is a tautology once the generic formulas are accepted: the degree-$D$ part of $V_D$ cannot exceed the number of degree-$D$ monomials, so a generic value above $\binom ND$ is impossible and the closure must have left the generic regime, i.e.\ produced additional degree falls. The ``if'' direction --- that \emph{every} excess, however small, triggers a cascade that fills the entire truncated ideal --- is an empirical fact: the smallest excesses observed to collapse are 8 out of 560 at $(16,16)$, 40 out of $46\,376$ at $(34,32)$ and 304 out of $201\,376$ at $(32,20)$, and no partial collapse was ever seen.

\subsection{Tests}

\begin{table}[ht]\centering\footnotesize
\caption{Boundary cases (both sides of a threshold, or margin below 3\%). c = collapse (degree $D$ suffices), g = generic (degree $D+1$ or more needed). ``M4GB'' cases were run with the corrected solver, ``closure'' cases by Lemma~\ref{lem:std}.}\label{tab:boundary}
\begin{tabular}{cp{3.9cm}p{4.3cm}p{1.4cm}p{1.6cm}p{2.0cm}}
\toprule
$D$ & $(N,n)$ & $Q_D-\binom ND$ & predicted & observed & method\\
\midrule
3 & (14,14), (15,15), (16,16) & $+63$, $+40$, $+8$ & c & c & closure\\
3 & (14,13) & $+26$ & c & c & closure\\
3 & (16,15), (17,17), (18,17), (18,18), (20,19) & $-35$, $-34$, $-136$, $-87$, $-285$ & g & g & closure\\
4 & (20,11), (20,12) & $-93$, $+285$ & g, c & g, c & closure\\
4 & (22,13), (22,14) & $-217$, $+266$ & g, c & g, c & closure\\
4 & (24,15), (24,16) & $-546$, $+54$ & g, c & g, c & closure\\
4 & (26,16), (28,16) & $-1918$, $-4899$ & g & g & closure\\
4 & (30,24), (30,25) & $-801$, $+195$ & g, c & g, c & closure and M4GB\\
4 & (28,21), (28,22) & $-504$, $+348$ & g, c & g, c & M4GB\\
4 & (32,28), (32,29) & $-302$, $+841$ & g, c & g, c & M4GB\\
4 & (34,31), (34,32) & $-1271$, $\mathbf{+40}$ & g, c & g, c & M4GB\\
4 & (34,33), (34,34), (36,36) & $+1342$, $+2635$, $+117$ & c & c & M4GB and closure\\
4 & (36,35), (38,37), (38,38), (40,39), (40,40) & $-1365$, $-5217$, $-3534$, $-10426$, $-8530$ & g & g & M4GB and closure\\
5 & 15 cases of Table~\ref{tab:level5} & from $-30\,816$ to $+3\,400$ & & all as predicted & closure\\
5 & (34,22), (34,23) & $-3058$, $+5932$ & g, c & g, c & M4GB\\
5 & (36,24), (36,25) & $-9936$, $\mathbf{+1308}$ & g, c & g, c & M4GB\\
5 & (40,30), (40,31) & $-7608$, $+8492$ & g, c & g, c & M4GB\\
\bottomrule
\end{tabular}
\end{table}

Together with the 55 grid systems of Section~\ref{sec:generic} (each of which is a test at $D=3$ and at $D=4$) the criterion was checked on 59 boundary cases and 110 grid outcomes without a single failure; 21 of the boundary predictions were recorded, with their margins, before the systems were run. The M4GB tests at $D=5$ are the most expensive computations of this note (up to 63\,h and 8.5\,GB at $(40,30)$) and are also the ones that reach closest to the square systems: for $(40,30)$ and $(40,31)$ the ratio $n/N$ is $0.75$--$0.78$, against $0.50$--$0.59$ for the points on which $Q_5$ was fitted.

\subsection{Consequences for the family of HKY}

For the square systems (even $n$, $N=n$) the criterion gives degree 3 up to $n=16$, degree 4 for $18\le n\le36$ and degree 5 for $38\le n\le52$; for odd $n$ ($N=n+1$) degree 4 from $n=15$, degree 5 from $n=35$ and degree 6 from $n=51$. The margins at the predicted transition are $-0.24\%$ at $(52,51)$, $+1.25\%$ at $(52,52)$, $-3.1\%$ at $(54,53)$ and $-1.6\%$ at $(54,54)$; the first of these is the only tight one. Every entry up to $n=40$ is confirmed by computation; the values for $n=41$--$54$ are predictions of the polynomial $Q_5$, extrapolated from $n/N\le0.78$ to $n/N=1$. A direct test would require a Gröbner-basis computation with 52 variables reaching step degree 6, which is beyond M4GB on a workstation by two to three orders of magnitude; it is the natural target for a faster solver.
\section{Structure of the closures and a growth law}\label{sec:structure}

\subsection{Semi-regular part and lag}

For $m$ Boolean quadratic polynomials in $N$ variables that are \emph{semi-regular} in the sense of Bardet, Faugère, Salvy and Yang \cite{BFS03,BFSY05}, the dimension of the degree-$k$ part of the Macaulay span is the coefficient
\[
I_k^{\SR}(N,m)=\binom Nk-[z^k]\,\frac{(1+z)^N}{(1+z^2)^m}
= m\binom N{k-2}-\binom{m+1}2\binom N{k-4}+\binom{m+2}3\binom N{k-6}-\cdots,
\]
as long as this is positive; the first $k$ for which the coefficient of the series becomes non-positive is the semi-regular degree of regularity. We verified this description on random systems: for 33 random dense Boolean quadrics in 22 variables the closure has the profile $(33,759,7821)$ predicted by the series, with no correction.

Our polynomials $Q_D$ are the semi-regular counts of a system of $m=3n$ quadrics, minus a lag:
\begin{align*}
Q_3&=I_3^{\SR}(N,3n)-\mathrm{lag}_3, & \mathrm{lag}_3&=5n+\tbinom N2,\\
Q_4&=I_4^{\SR}(N,3n)-\mathrm{lag}_4, & \mathrm{lag}_4&=87n,\\
Q_5&=I_5^{\SR}(N,3n)-\mathrm{lag}_5, & \mathrm{lag}_5&=n(79N-660).
\end{align*}
The identification of the Koszul-type term $\binom{3n+1}2$ in $Q_4$ and of $\binom{3n+1}2N$ in $Q_5$ with the second term of the semi-regular series is what makes this decomposition natural rather than a rewriting: the $3n$ quadrics are exactly the polynomials of $W_2$ (Section~\ref{sec:generic}), and the closures behave as if these $3n$ polynomials were in generic position, up to a lag that is small relative to the leading term ($26\%$, $5.0\%$ and $6.4\%$ of $I_D^{\SR}$ at the three thresholds). The lag is not a property of the ideal: the truncations $J_{\le k}$ of the ideal itself are fixed by the solution set ($|\B_{\le k}|$ minus the number of standard monomials of degree $\le k$) and carry no generic information. It is a property of what the degree-capped closure can reach.

\subsection{The first stage exactly}

The closure has two stages: the span $R_D$ of the plain products $\{t\,q: q\in W_2,\ \deg t\le D-2\}$ and the products of variables with the degree falls that appear inside that span. The first stage admits an exact block decomposition. Write $\ell$ for the trace-linear polynomial, $m'=3n-N$ for the number of quadrics of $W_2$ not in the ideal of $\ell$, and $M_3=2n-N+1$ for the number of degree falls at level 3 (``mutants'' in the terminology of \cite{Ding08,BDMM09}). Then
\[
R_D=\binom{N-1}{\le D-1}+\SR_{\le D}(N-1,m')-M_3\binom{N}{\le D-3}-\delta_D,
\]
where $\binom{N-1}{\le D-1}=\dim\ell\,\B_{\le D-1}$ is the contribution of the ideal of $\ell$, where $\SR_{\le D}(N-1,m')=\sum_{k=2}^D I_k^{\SR}(N-1,m')$ is the semi-regular count of the remaining quadrics in the $N-1$ variables of the quotient by $\ell$, the third term counts the defining relations of the mutants multiplied by monomials, and
\[
\delta_3=\delta_4=0,\qquad \delta_5=\binom N2-1-3\binom{M_3}2-\frac{(M_3-1)(N+1)}2 .
\]
At $D=4$ the formula reproduces the rank of the product set on all 7 systems where it was computed ($2\,994$ at $(20,10)$; $4\,025$, $4\,636$, $5\,238$ at $N=22$; $5\,269$, $6\,008$, $6\,738$ at $N=24$); at $D=5$, with $\delta_5$ fitted on $(24,12)$, $(30,15$--$17)$ and $(32,16)$, it reproduces $R_5$ on 9 systems, including the predicted values $46\,372$, $63\,386$ and $71\,178$ at $(26,13)$, $(28,14)$, $(28,15)$ ($M_3=1,1,3$). At $D=3$ the formula is an identity: $\binom{N-1}2+m'(N-1)-M_3=3nN-5n-\binom N2=Q_3$, so $Q_3$ is \emph{derived} from the block structure. The new term $\delta_5$ has an evident reading --- $\binom N2-1$ relations between the ideal of $\ell$ and the other quadrics that first appear at degree 5, three Koszul-type relations per pair of mutants, and $(M_3-1)(N+1)/2$ relations between mutants and $\ell$ --- but we have not proved it.

The second stage, $S_D=P_D-R_D$, is known exactly as a difference of the formulas above ($S_4=N^3/3-N^2/2+2Nn-5N/6-85n+1$; $S_5$ of degree 4) and is the part of the closure that no Macaulay-type count captures: at $(20,10)$ it contributes $2\,001$ of the $4\,365$ degree-4 dimensions. Its structural derivation is open (Section~\ref{sec:open}).

\subsection{Thresholds and the growth of $\Dreg$}

For square systems ($n=N$) the criterion $Q_D(N,N)>\binom ND$ gives the last $N$ at which degree $D$ suffices. Table~\ref{tab:thresholds} compares the semi-regular thresholds (lag $=0$) with the actual ones.

\begin{table}[ht]\centering\small
\caption{Last $N$ with degree $D$ for square systems.}\label{tab:thresholds}
\begin{tabular}{ccccc}
\toprule
$D$ & semi-regular, $m=3N$ & with lag (this note) & ratio & lag$/I^{\SR}$ at threshold\\
\midrule
3 & 20 & 16 & 0.80 & 26\%\\
4 & 37 & 36 & 0.97 & 5.0\%\\
5 & 56 & 52 (predicted) & 0.93 & 6.4\%\\
6 & 76 & 66--76 & & \\
7 & 97 & & & \\
8 & 119 & & & \\
\bottomrule
\end{tabular}
\end{table}

The semi-regular thresholds grow with gaps $17,19,20,21,22,22,22,23,\dots$, tending to the asymptotic BFSY slope: for $m=\alpha N$ Boolean quadrics the semi-regular degree of regularity is $D\approx c(\alpha)N$ with
\[
c(3)=-3+\tfrac12+\tfrac12\sqrt{2\cdot9-30-1+2\cdot5\sqrt{15}}\approx0.036,
\]
i.e.\ one degree per 28 variables asymptotically, one per 20 in the computed range. The actual thresholds follow the semi-regular ones at $93$--$97\%$ for $D\ge4$. Unless the lag grows unexpectedly, the degree of regularity of the Weil-descent system of $S_3$ therefore grows \emph{linearly} in $n$ --- about 6 for $53\le n\lesssim70$, 7 to about 90, 8 to about 110 --- and not as a constant, as the first fall degree assumption would require, nor logarithmically. The mechanism is the multiplier 3: the ideal contains three times as many quadrics as there are equations, so the system is as easy as a semi-regular system of $3n$ quadrics (for $m=n$ the slope would be $c(1)\approx0.09$), but no easier.

For the index calculus this changes nothing qualitatively: the decomposition with two points is exponential in $n$ with a smaller constant than the generic estimate, and the relation collection needs $2^{n/2}$ decompositions regardless.

\section{Three-point decompositions: the Weil descent of $S_4$}\label{sec:s4}

\subsection{The systems}

The subexponential estimate of \cite{PQ12} uses decompositions into $m\ge3$ points, so the assumption that matters most is the one for $S_4$ and beyond. With $S_3(X_1,X_2,X)=AX^2+BX+C$, $A=(X_1+X_2)^2$, $B=X_1X_2$, $C=B^2+a_6$, the fourth summation polynomial is the resultant of two copies of $S_3$ in the common variable, which in characteristic 2 reads
\[
S_4(X_1,X_2,X_3,X_4)=(AC'+A'C)^2+(AB'+A'B)(BC'+B'C),
\]
primes denoting the same expressions in $(X_3,X_4)$. We checked on random points that it vanishes on $x$-coordinates of $P_1+P_2+P_3=Q$ and not on random quadruples. Substituting $x(P_i)=\sum_j u_{ij}v_j$ with $u_{ij}\in\F_2$ and $x(P_4)=x(Q)$, and using that squaring is $\F_2$-linear, the descended system consists of $n$ equations in $N=3d$ variables whose degree over $\F_2$ is exactly 6 (a monomial $X_i^3$ contributes degree 2, $X_i,X_i^2,X_i^4$ degree 1); every equation has about $d^6/2$ monomials. We take $d=\lceil n/3\rceil$ for square systems and smaller $n$ for the same $d$ to move along the $(N,n)$ family, as for $S_3$. Solutions of the descended system are counted exhaustively over $V_1\times V_2\times V_3$ (see the remark in Section~8.2), and every descended system was verified to vanish on all of them.

Two consequences of degree 6 shape the computation. At $D=6$ nothing happens: there are no products to form. At $D=7$ the products $x_af$ of the equations already fall in degree, so the first fall degree is $7=\deg+1$, and the question is whether $V_7$ contains the Gröbner basis. The monomial counts are far larger than for $S_3$ ($|\B_{\le7}|=536\,155$ at $N=24$), and M4GB, whose strength lies in keeping tails fully reduced, builds matrices of over $2\cdot10^6$ rows already at $N=21$ and stops with \texttt{std::bad\_alloc}; the third closure implementation of Section~\ref{sec:tools}, with the echelon in fully reduced form, is what made $N=24$ possible.

\subsection{Results}

\begin{table}[ht]\centering\small
\caption{$S_4$ at $D=7$: in every case $\dim V_7=|\B_{\le7}|-\#\text{sol}$ and all standard monomials have degree $\le6$, so the reduced Gröbner basis lies in $V_7$ and $\Dsolve=7$. \#sol is the number of $\F_2$-solutions of the descended system (see the remark below). For the four systems marked $\dagger$ the closure at $D=8$ was computed as well and has the same standard monomials.}\label{tab:s4}
\begin{tabular}{rrrrl}
\toprule
$N$ & $n$ & $n/N$ & \#sol & $\dim V_7$\\
\midrule
12 & 12 & 1 & 2 & 3\,300\\
15 & 15 & 1 & 3 & 16\,381\\
18 & 18 & 1 & 1 & 63\,003\\
18 & 16, 14, 12, 10, 9 & 0.89--0.50 & 2, 18, 71, 363, 809 & $|\B_{\le7}|-\#\text{sol}$\\
18 & 8 & 0.44 & 842 & 62\,162$^\dagger$\\
18 & 7 & 0.39 & 1\,791 & 61\,213$^\dagger$\\
18 & 6 & 0.33 & 4\,933 & 58\,071$^\dagger$\\
21 & 21 & 1 & 2 & 198\,438\\
21 & 18, 15, 12 & 0.86--0.57 & 5, 62, 499 & $|\B_{\le7}|-\#\text{sol}$\\
21 & 9 & 0.43 & 4\,538 & 193\,902$^\dagger$\\
24 & 24 & 1 & 1 & 536\,154\\
24 & 12 & 0.50 & 3\,881 & 532\,274\\
\bottomrule
\end{tabular}
\end{table}

Every one of the twenty systems collapsed at $D=7$: the closure contains the whole truncated ideal, $\Dsolve$ equals the first fall degree $7$, and this holds from square systems down to $n=N/3$, i.e.\ for systems with thousands of solutions. No generic regime was observed at all --- for $S_3$ it sets in at $N=17$ for $D=3$ and the collapse boundary $n_0(N)$ then moves by about $0.025N$ per variable at fixed $D$ (Section~\ref{sec:criterion}); for $S_4$ up to $N=24$ there is no boundary to move. The closures at $D=8$ for the four systems with the most solutions add no standard monomial and confirm the picture.

\paragraph{Remark (counting solutions).} Lemma~\ref{lem:std} needs the number of $\F_2$-solutions of the \emph{descended} system, and for $S_4$ this is not the number of decompositions. The resultant $S_4$ vanishes whenever the two quadratics have a common root in the algebraic closure, so also when the common root lies in $\F_{2^{2n}}$ --- these solutions have no point of $E(\F_{2^n})$ behind them --- and when both leading coefficients vanish ($x_1=x_2$, $x_3=x(Q)$). A count through common roots in $\F_{2^n}$ misses them: at $(18,8)$ it gives 833 where the system has 842 solutions, and the nine ``missing'' elements of $V_7$ that we first took for a genuine deficit were exactly these. All counts in Table~\ref{tab:s4} are exhaustive over $V^3$; the descended systems were verified to vanish on every solution found.

Within the computed range the first fall degree assumption therefore holds for $m=3$ while it fails for $m=2$. Whether degree 7 keeps sufficing for larger $N$ is open; the next value, $N=27$, needs about $37$\,GB for the closure and is the natural target for a larger machine.

\section{Controls}

\paragraph{Random quadrics.} As noted in Section~\ref{sec:structure}, $3n$ random dense Boolean quadrics reproduce the semi-regular series exactly (no lag). Random \emph{bilinear} quadrics in the same two blocks of variables collapse earlier (fewer monomials are available to bilinear products); the Weil-descent systems, although bilinear in their quadratic part, behave like dense ones because $W_2$ is not bilinear.

\paragraph{Subfield subspaces.} If $n=2k$ and $V=\F_{2^k}$ is the subfield --- the classical setting of Gaudry and Diem for extension degree 2 --- the span of the $n$ equations contains $k$ linear polynomials instead of one, $\dim V_2$ rises from $2n-1=31$ to 99 at $n=16$, and $\Dsolve=3$ for every $n$ we tried (16 to 24). This is not a failure of the model but its confirmation on a degenerate case: for $X_0,X_1\in\F_{2^k}$ the equation $S_3=0$ splits over $\F_{2^k}$ into $Z=pS+a_{6,1}$ and $p^2S^2+qS+a_{6,0}=0$ with $S=X_0^2+X_1^2$, $Z=X_0X_1$ and $x(Q)^2=p\,x(Q)+q$, the second being Frobenius-linear in $S$; the decomposition is a linear system followed by one quadratic equation. HKY's random subspaces avoid exactly this structure.

\paragraph{Instance independence.} All formulas of Sections~\ref{sec:generic}--\ref{sec:structure} were checked on a second instance at $(28,15)$ and $(28,16)$ (level 5) and on 378 instances at $n\le24$ (levels 3--4): the dimensions agree to the unit.

\section{Open questions}\label{sec:open}

\begin{enumerate}
\item \emph{Why $3n$ and $2n-N+1$.} The dimension of the quadratic part of the ideal and the number of mutants were measured, not derived. Both should follow from the $\F_{2^n}$-linear structure of $S_3$: the equations are the coordinates of a single polynomial over $\F_{2^n}$, and the mutants are combinations $\Tr(\gamma\Lambda S_3)$ with $\Lambda$ ranging over $\F_{2^n}$-linear forms in the coordinates, whose cubic part vanishes; the count $2n-N+1$ has the shape of a dimension count ($2n$ unknowns, $N$ constraints, one degree of freedom) that we could not complete.
\item \emph{The second stage.} A closed form for the rank $S_D$ of the products of variables with the degree falls would give $\mathrm{lag}_D$ for every $D$ and turn the $D=6$ threshold, currently bracketed as $66\le N\le76$, into a prediction. The data for $S_4$ (exact) and $S_5$ (nine points) are in the archive.
\item \emph{Semi-regularity of $W_2$.} Is the Macaulay span of $W_2$ semi-regular in the quotient by $\ell$, as the first-stage formula asserts, for all $D$? (For the complexity of Weil restriction systems in general see \cite{HKYY18,CCG23}.) A proof would settle the first stage; the deviation $\delta_5$ shows that the naive statement needs the corrections listed in Section~\ref{sec:structure}.
\item \emph{Beyond $N=24$ for three points.} Section~\ref{sec:s4} leaves the central question open: does degree 7 keep sufficing for the Weil descent of $S_4$ --- in which case the first fall degree assumption holds for $m=3$ and the estimate of \cite{PQ12} stands for that step --- or does a generic regime appear at some $N>24$, as it does at $N=17$ for $S_3$? One more value of $N$ ($27$, about $37$\,GB) is the next test; an argument for why the degree-7 closure of an $S_4$ system reaches the whole truncated ideal would be better still.
\end{enumerate}

\section{Reproducibility}

All systems are generated from seeds recorded in the supplementary archive (\url{https://github.com/alekssolov/weil-descent-s3-dreg}); the generator (Python~3, standard library and numpy), the Python closure and its verification by Buchberger--Möller, the C closure program with its Macaulay mode, the Buchberger reference implementation with switchable pair criteria, the emulation of M4GB's criteria, the M4GB patch (as a script and as a unified diff), the $S_4$ generator with its self-checks, the memory-light closure program, every input file run in M4GB and every solver log are provided. The $S_4$ closures at $N=24$ took about 72 hours each and 10\,GB on an 8-thread desktop (two of them run concurrently, hence partly in swap); the degree-5 closures of the $S_3$ systems with 40 variables took 32 hours each with a peak of 20\,GB. Timings: level-5 closures at $N=32$ take 5--13\,h and 3.4--3.7\,GB on eight threads; the largest solver runs are $(40,30)$ --- 63\,h and 8.5\,GB --- and the square system $n=40$ --- 6\,h and 1.5\,GB. Two practical remarks for anyone repeating the solver runs are in Section~\ref{sec:m4gb}: the libtool wrapper in \texttt{bin/} is not the executable, and the pair criteria of releases before 25 September 2026 need the correction of pull request \#14 for step degrees to be meaningful.

\paragraph{Acknowledgements.} We are deeply grateful to Dr Marc Stevens and Dr Rusydi H. Makarim for making their Gröbner-basis solver M4GB publicly available; all Gröbner-basis computations in this note were carried out with it, and we thank Dr Stevens for his prompt response to our report on the pair criteria and for merging the correction. The author used an AI assistant (Anthropic Claude) for programming, data analysis and drafting; all computations were run and all results verified by the author.

\begin{thebibliography}{99}\small
\bibitem{ACFP12} M.~R. Albrecht, C.~Cid, J.-C. Faugère, L.~Perret, On the relation between the MXL family of algorithms and Gröbner basis algorithms, J.~Symbolic Comput.\ 47 (2012), 926--941.
\bibitem{BFS03} M.~Bardet, J.-C. Faugère, B.~Salvy, Complexity of Gröbner basis computation for semi-regular overdetermined sequences over GF(2) with solutions in GF(2), Research Report RR-5049, INRIA, 2003.
\bibitem{BFSY05} M.~Bardet, J.-C. Faugère, B.~Salvy, B.-Y. Yang, Asymptotic behaviour of the degree of regularity of semi-regular polynomial systems, in: MEGA 2005, Eighth International Symposium on Effective Methods in Algebraic Geometry, Porto Conte, Italy, 2005.
\bibitem{Magma} W.~Bosma, J.~Cannon, C.~Playoust, The Magma algebra system I: the user language, J.~Symbolic Comput.\ 24 (1997), 235--265.
\bibitem{CCG23} A.~Caminata, M.~Ceria, E.~Gorla, The complexity of solving Weil restriction systems, J.~Algebra 621 (2023), 116--133.
\bibitem{CG23} A.~Caminata, E.~Gorla, Solving degree, last fall degree, and related invariants, J.~Symbolic Comput.\ 114 (2023), 322--335.
\bibitem{Diem11} C.~Diem, On the discrete logarithm problem in elliptic curves, Compositio Math.\ 147 (2011), 75--104.
\bibitem{Diem13} C.~Diem, On the discrete logarithm problem in elliptic curves II, Algebra \& Number Theory 7 (2013), 1281--1323.
\bibitem{BDMM09} J.~Buchmann, J.~Ding, M.~S.~E. Mohamed, W.~S.~A.~E. Mohamed, MutantXL: solving multivariate polynomial equations for cryptanalysis, Dagstuhl Seminar Proceedings 09031 (Symmetric Cryptography), 2009.
\bibitem{Ding08} J.~Ding, J.~Buchmann, M.~S.~E. Mohamed, W.~S.~A.~E. Mohamed, R.-P. Weinmann, MutantXL, in: Proceedings of the 1st International Conference on Symbolic Computation and Cryptography (SCC 2008), Beijing, LMIB, 2008, 16--22.
\bibitem{F4} J.-C. Faugère, A new efficient algorithm for computing Gröbner bases (F4), J.~Pure Appl.\ Algebra 139 (1999), 61--88.
\bibitem{FHJRV14} J.-C. Faugère, L.~Huot, A.~Joux, G.~Renault, V.~Vitse, Symmetrized summation polynomials: using small order torsion points to speed up elliptic curve index calculus, EUROCRYPT 2014, LNCS 8441, 40--57.
\bibitem{FW18} J.-C. Faugère, A.~Wallet, The point decomposition problem over hyperelliptic curves: toward efficient computation of discrete logarithms in even characteristic, Des.\ Codes Cryptogr.\ 86 (2018), 2279--2314.
\bibitem{Gaudry09} P.~Gaudry, Index calculus for abelian varieties of small dimension and the elliptic curve discrete logarithm problem, J.~Symbolic Comput.\ 44 (2009), 1690--1702.
\bibitem{GM88} R.~Gebauer, H.~M. Möller, On an installation of Buchberger's algorithm, J.~Symbolic Comput.\ 6 (1988), 275--286.
\bibitem{Huang23} M.-D. A. Huang, Last fall degree of semi-local polynomial systems, arXiv:2311.02804 (2023).
\bibitem{HKY15} M.-D. A. Huang, M.~Kosters, S.~L. Yeo, Last fall degree, HFE, and Weil descent attacks on ECDLP, CRYPTO 2015, LNCS 9215, 581--600; IACR ePrint 2015/573.
\bibitem{HKYY18} M.-D. A. Huang, M.~Kosters, Y.~Yang, S.~L. Yeo, On the last fall degree of zero-dimensional Weil descent systems, J.~Symbolic Comput.\ 87 (2018), 207--226; arXiv:1505.02532.
\bibitem{JV12} A.~Joux, V.~Vitse, Cover and decomposition index calculus on elliptic curves made practical, EUROCRYPT 2012, LNCS 7237, 9--26.
\bibitem{KY15} M.~Kosters, S.~L. Yeo, Notes on summation polynomials, arXiv:1503.08001 (2015).
\bibitem{MS17} R.~H. Makarim, M.~Stevens, M4GB: an efficient Gröbner-basis algorithm, ISSAC 2017, 293--300; source code \url{https://github.com/cr-marcstevens/m4gb}.
\bibitem{Nagao10} K.~Nagao, On the complexity of the decomposition attack, IACR ePrint 2010/511.
\bibitem{PQ12} C.~Petit, J.-J. Quisquater, On polynomial systems arising from a Weil descent, ASIACRYPT 2012, LNCS 7658, 451--466.
\bibitem{Semaev04} I.~Semaev, Summation polynomials and the discrete logarithm problem on elliptic curves, IACR ePrint 2004/031.
\end{thebibliography}

\end{document}
